% 型3：見開き対応型
% 左ページに問題，右ページに解答を置き，一つの見開きで完結させる．
% ページをめくらずに答え合わせができるため，短時間の自習に向く．
% 1段組にして，左右の対応を目で追えるようにする．
\PassOptionsToPackage{twocolumn=false}{surikumi}
\documentclass[lualatex,paper=b5j,fontsize=10pt,twoside]{jlreq}

\usepackage{surikumi-mondaishu}

\MathMagazineSetup{
  feature={},
  series={見開き演習／数学 I},
  series-position=left,
  pattern=checker,
  title={二次関数の最大・最小},
  author={入試基礎講座},
  author-font=mincho,
  author-position=right,
  title-fill=black,
  deck={定義域が動く場合の場合分けを，見開き一組で確認する．左が問題，右が解答である．}
}

\begin{document}

\MakeMathMagazineTitle

この見開きでは，軸と定義域の位置関係で場合分けをする型を扱う．
左ページの 4 問を解いてから，右ページで確認する．

\MSStartVerso

\MSSpreadHeading{問題}

\begin{mmproblems}
  \item $f(x)=x^2-4x+1$ の $0\leq x\leq 1$ における最小値を求めよ．

  \item $f(x)=x^2-4x+1$ の $0\leq x\leq 3$ における最小値を求めよ．

  \item $a$ を正の定数とする．$f(x)=x^2-4x+1$ の $0\leq x\leq a$ における
    最小値を $a$ で表せ．

  \item $a$ を実数とする．$f(x)=x^2-2ax+1$ の $0\leq x\leq 1$ における
    最小値を $a$ で表せ．

  \item $f(x)=x^2-4x+1$ の $0\leq x\leq 1$ における最大値を求めよ．

  \item $f(x)=-x^2+2x+3$ の $0\leq x\leq 3$ における最大値と最小値を求めよ．

  \item $a$ を実数とする．$f(x)=x^2-2ax+1$ の $0\leq x\leq 1$ における
    最大値を $a$ で表せ．

  \item $a$ を実数とする．$f(x)=x^2-2x+3$ の $a\leq x\leq a+1$ における
    最小値を $a$ で表せ．
\end{mmproblems}

\clearpage

\MSSpreadHeading{解答}

\begin{mmproblems}
  \item $f(x)=(x-2)^2-3$ であるから軸は $x=2$ である．
    定義域 $0\leq x\leq 1$ は軸の左側にあり，$f$ はこの範囲で減少する．
    よって $x=1$ で最小となり，最小値は
    \[
      f(1)=-2
    \]

  \item 軸 $x=2$ は定義域 $0\leq x\leq 3$ に含まれる．
    よって頂点で最小となり，最小値は
    \[
      f(2)=-3
    \]

  \item 軸 $x=2$ が定義域 $0\leq x\leq a$ に入るかどうかで分かれる．

    \MMCaseHeading{1}{$0<a<2$ のとき}
    定義域は軸の左側にあり，$f$ は減少する．よって $x=a$ で最小となり，
    最小値は $a^2-4a+1$ である．

    \MMCaseHeading{2}{$a\geq 2$ のとき}
    軸が定義域に含まれるから，最小値は $f(2)=-3$ である．

  \item $f(x)=(x-a)^2+1-a^2$ であるから軸は $x=a$ である．

    \MMCaseHeading{1}{$a<0$ のとき}
    定義域は軸の右側にあり，$f$ は増加する．よって $x=0$ で最小となり，
    最小値は $1$ である．

    \MMCaseHeading{2}{$0\leq a\leq 1$ のとき}
    軸が定義域に含まれるから，最小値は $1-a^2$ である．

    \MMCaseHeading{3}{$a>1$ のとき}
    定義域は軸の左側にあり，$f$ は減少する．よって $x=1$ で最小となり，
    最小値は $2-2a$ である．

  \item 定義域 $0\leq x\leq 1$ は軸 $x=2$ の左側にあり，$f$ は減少する．
    よって $x=0$ で最大となり，最大値は
    \[
      f(0)=1
    \]

  \item $f(x)=-(x-1)^2+4$ であるから軸は $x=1$ で，グラフは上に凸である．
    軸は定義域 $0\leq x\leq 3$ に含まれるから，最大値は $f(1)=4$ である．
    最小値は軸から遠いほうの端点でとる．$f(0)=3$，$f(3)=0$ であるから，
    最大値と最小値は
    \[
      4,\quad 0
    \]

  \item 下に凸のグラフでは，最大値は軸から遠いほうの端点でとる．
    定義域 $0\leq x\leq 1$ の中点は $x=\dfrac{1}{2}$ である．

    \MMCaseHeading{1}{$a<\dfrac{1}{2}$ のとき}
    軸は中点より左にあり，$x=1$ のほうが遠い．最大値は $f(1)=2-2a$ である．

    \MMCaseHeading{2}{$a\geq\dfrac{1}{2}$ のとき}
    軸は中点以上にあり，$x=0$ のほうが遠い．最大値は $f(0)=1$ である．

  \item $f(x)=(x-1)^2+2$ であるから軸は $x=1$ である．
    今度は定義域 $a\leq x\leq a+1$ のほうが動く．

    \MMCaseHeading{1}{$a+1<1$ すなわち $a<0$ のとき}
    定義域は軸の左側にあり，$f$ は減少する．よって $x=a+1$ で最小となり，
    最小値は $a^2+2$ である．

    \MMCaseHeading{2}{$0\leq a\leq 1$ のとき}
    軸が定義域に含まれるから，最小値は $f(1)=2$ である．

    \MMCaseHeading{3}{$a>1$ のとき}
    定義域は軸の右側にあり，$f$ は増加する．よって $x=a$ で最小となり，
    最小値は $a^2-2a+3$ である．
\end{mmproblems}

\end{document}
